\scanpage{17}
1) $f:X\longrightarrow S$ projectif, $X$ régulier,
$S$ \unsure{spectr.|spectre,|spectre} ann. de val. discrète
pt générique $\eta$, pt spécial $s$

\[
\left\{\begin{aligned}
&X_\eta \text{ est lisse}\\
&X_s \text{ est géom\textsuperscript{t} intègre }\deleted{\unsure{\text{algébr.}}}
\end{aligned}\right.
\]

\[
 \unsure{\Longleftrightarrow|\Longrightarrow|\text{\deleted{Soit}}}\quad\dim X_\eta=n\qquad\dim X=n+1
\]

$C,D$ deux cycles sur $X$, de codimensions
$c$ et $d$ avec $c+d=n+1$. On suppose

a) \deleted{\illegible{1}} $C_\eta$ et $D_\eta$ sont alg. équiv. à $0$
\deleted{\unsure{et} et $\operatorname{supp}C_\eta\cap\operatorname{supp}D_\eta=\varnothing$}

b) $C$ et $D$ \unsure{n’aient|n’ont} que des \unsure{composants|composantes} \unsure{horizontales|horizontaux|horizontal},
i.e. $C=\overline{C_\eta}$, $D=\overline{D_\eta}$.

c) $C.D$ est défini $[\Longrightarrow\operatorname{supp}C_\eta\cap\operatorname{supp}D_\eta=\varnothing]$

Considérons $i(C.D)$, il ne dépend que de
$C_\eta$ et $D_\eta$, cycles \deleted{alg.} de $X_\eta$ \inserted{alg. équ. à $0$} de codimensions
$c,d$ tels que $c+d=n+1$, \deleted{On veut} \inserted{de supports}
disjoints. On veut alors donner une
description de $i(C.D)$ en termes
de $C_\eta,D_\eta$, comme un

\[
 \langle C_\eta,D_\eta\rangle_v
\]

\uline{indépendant du modèle $X$} \inserted{\uline{ou de son existence}},
(\uline{à valeurs dans $\mathbb Q$} \inserted{à priori})
% archive sheet number: 16
